% !PRINTED 512
\weberpage{512}

%<<CONT
\gesperrt{odd divisor, and in such a way that, if $x$ is even, $y\equiv 2\pmod 4$, and further in such a way that
\[
n=\tfrac14(y^2-Dx^2)
\]
is an odd square number. Let $M$ further be determined by formula (12). Then the function $H_{-D}(u)$ has a divisor
\weq{22}{\Phi(M,u),}
which with none of the functions
\weq{23}{\Phi(-M,u),\quad \Phi(iM,u),\quad \Phi(-i,Mu)}
% [sic] German source prints $\Phi(-i,Mu)$, misplacing the comma; the sign orbit requires $\Phi(-iM,u)$
has a root in common.}

\begin{center}
§ 137.\quad \gesperrt{Reduction of odd discriminants to even ones.}
\end{center}

According to a theorem stated by \gesperrt{Kronecker}, the class equation can, upon the adjunction of square roots, be decomposed into as many factors as there are genera of the classes of forms for the class equation in question, and indeed in such a way that each of these factors is satisfied only by those class invariants for which the corresponding quadratic forms $(A,B,C)$ belong to one and the same genus. We obtain this decomposition from Theorem 1, proved in the preceding section. % [sic] German source prints the ungrammatical "aus den ... Satz"

That this theorem, in the form in which we have stated it, refers only to even discriminants $D$ is no restriction for the proof of \gesperrt{Kronecker}'s theorem.

For if $D\equiv 1\pmod 8$, then by § 123 the class field $\mathfrak{K}(D)$ is identical with the class field $\mathfrak{K}(4D)$. But if $D\equiv 5\pmod 8$, then the degree of $\mathfrak{K}(4D)$ is three times as great as that of $\mathfrak{K}(D)$; the latter, however, is contained in the former, and the number of genera for $\mathfrak{K}(D)$ and $\mathfrak{K}(4D)$ is the same (§ 104).

If $u=j(\omega)$ is a class invariant for $D\equiv 1\pmod 8$, then
\weq{1}{v=j(\omega')=j(2\omega),\quad j\left(\dfrac{\omega}{2}\right),\quad j\left(\dfrac{\omega+1}{2}\right)}
are class invariants of $4D$, and $u$ is a rational function of $v$:
\weq{2}{u=\varphi(v),}
which remains unchanged, whichever of the three quantities (2) one may put for $v$. % [sic] German source cites (2), but the three quantities occur in (1)
All characters of the forms whose roots are $\omega$ and $\omega'$ have the same value.
